As large language models (LLMs) are increasingly deployed as evaluators and judges, their sensitivity to user tone becomes a critical concern. While prior work highlights the risk of sycophancy when models face user pushback, we investigate a parallel phenomenon: social facilitation. Specifically, we explore whether exposing an LLM to skeptical or judgmental tones can act as a catalyst for deeper, more rigorous reasoning. In this work, we present a study on the ``Judgment Effect,'' examining how varying prompt tones---from neutral to highly judgmental---affect the length and quality of reasoning traces in \gptmini. We evaluate the model on two reasoning tasks: social arbitration using the \raio dataset and causal auditing via \causalt. Our results show a clear monotonic increase in reasoning effort under evaluative pressure, with the \judgmental condition extending trace lengths by up to 22.0\%. Crucially, on the complex social judgments of \raio, this heightened reasoning effort yields a 6.6\% improvement in accuracy over the neutral baseline, with minimal sycophantic label flipping. These findings indicate that controlled, judgmental framing can successfully induce deeper logical exploration and edge-case analysis in LLMs, providing a practical prompting strategy for high-stakes evaluation tasks.